\newif\ifglobal

%%%%%%%%%%%%%%%%%%%%%%%
%%% Compile to view %%%
%%%%%%%%%%%%%%%%%%%%%%%

%%%%%%%%%%%%%%%%%%%%%%%%%%%%%%%%%%%%%%%%%%%%%%%%%%%
%%% Readme:                                     %%%
%%% https://www.overleaf.com/read/bwdjvfjksvjy/ %%%
%%%%%%%%%%%%%%%%%%%%%%%%%%%%%%%%%%%%%%%%%%%%%%%%%%%

% >>> M?: marks a module setting number or important notice

% >>> M4: Set {./relative-path} to external files for this module <<<
\def \modulefiles {./01-RISCVCPU}
% To add an external file, use:
% (Replace '---' with actual filename)
% '\input{\modulefiles/tables/---.tex}' for tables
% '\includegraphics{\modulefiles/figures/---}' for figures
% '\subfile{\modulefiles/---}' for register files


% >>> M1: This module is in global project? <<<
% 'YES' - uncomment; 'NO' - comment out
\globaltrue

\ifglobal
    \documentclass[main\_\_EN.tex]{subfiles}

\else
    \documentclass[openany, 10pt]{book}

    % >>> M2: In ./00-shared/preamble-trm-module, also find
    %         settings for visibility of todo notes, watermark <<<
    \usepackage{preamble-trm-module}

    % >>> M3: Temporary labels in English,
    %         you can add your own labels there <<<
    \myexternaldocument{./00-shared/temp-labels-en}

\fi %\ifglobal


\setcounter{secnumdepth}{4}


% >>> M5: If this module needs any updates in
% './00-shared/preamble-trm-module.sty'
% (include additional package, etc.), contact your module PM
% or create the file './00-shared/changelog.tex' make the
% necessary updates yourself and document them in this file <<<


\begin{document}
% Sets language into which pre-defined bits of text are translated
\selectlanguage{english}

% Do not delete this block
\ifglobal
% Add the following front matter if
% this module is outside of global project
\else
    \InputIfFileExists{readme}{}{}
    {\let\clearpage\relax \listoftodos}
    {\let\clearpage\relax \listofdonetodos}
    \tableofcontents
    %\listoftables
    %\listoffigures
\fi

%%%%%%%%%%%%%%%%%%%%%%%%%
%%% TRM Module Begins %%%
%%%%%%%%%%%%%%%%%%%%%%%%%

\hypertarget{riscvcpu}{}
\chapter{ESP-RISC-V CPU}
\label{mod:riscvcpu}



\section{Overview}

ESP-RISC-V CPU is a 32-bit core based upon RISC-V ISA comprising base integer (I), multiplication/division (M) and compressed (C) standard extensions. The core has 4-stage, in-order, scalar pipeline optimized for area, power and performance. CPU core complex has an interrupt-controller (INTC), debug module (DM) and system bus (SYS BUS) interfaces for memory and peripheral access.
\begin{figure}[h]
\centering
\includegraphics[scale=1,width=0.8\textwidth]{\modulefiles/figures/CPU_C3_BLK_DIAG_v0.1}
\caption{CPU Block Diagram}
\end{figure}


\section{Features}
\label{sec:riscvcpu-features}

The ESP-RISC-V CPU has the following features:

\begin{itemize}

    \item Operating clock frequency up to 120 MHz
    \item Zero wait cycle access to on-chip SRAM and Cache for program and data access over IRAM/DRAM interface
    \item Interrupt controller (INTC) with up to 31 vectored interrupts with programmable priority and threshold levels
    \item Debug module (DM) compliant with RISC-V debug specification v0.13 with external debugger support over an industry-standard JTAG/USB port
    \item Debugger direct system bus access (SBA) to memory and peripherals
    \item Hardware trigger compliant to RISC-V debug specification v0.13 with up to 2 breakpoints/watchpoints
    \item Physical memory protection (PMP) for up to 16 regions
    \item 32-bit AHB system bus for peripheral access
    \item Configurable events for core performance metrics

\end{itemize}


\section{Address Map}

Below table shows address map of various regions accessible by CPU for instruction, data, system bus peripheral and debug.

\begin{longtable}{| l| l |c | c | c  |}
\caption{CPU Address Map} \label{tab:riscvcpu-mem-blocks}
\\ \hline
\rowcolor{lightgray}
Name & Description    &  Starting Address  &  Ending Address  & Access   \\ \hline
IRAM\label{TBD1} & Instruction Address Map  &   0x4000\_0000  &  0x47FF\_FFFF  &  R/W  \\ \hline
DRAM\label{TBD2} & Data Address Map &     0x3800\_0000  &  0x3FFF\_FFFF  &  R/W  \\ \hline
DM \label{TBD4} & Debug Address Map  &      0x2000\_0000  &  0x27FF\_FFFF  &  R/W  \\ \hline
AHB\label{TBD3} & AHB Address Map &      *default  &  *default  &  R/W  \\ \hline
\end{longtable}

*default : Address not matching any of the specified ranges (IRAM, DRAM, DM) are accessed using AHB bus.

%%%%%%%%%%%%%%%%%%%%%%%%%%%
%%% CS Register Summary %%%
%%%%%%%%%%%%%%%%%%%%%%%%%%%


\section{Configuration and Status Registers (CSRs)}
\hypertarget{riscvcpu-cpu-reg-summ}{}
\subsection{Register Summary}
\label{sec:riscvcpu-cpu-reg-summ}


Below is a list of CSRs available to the CPU. Except for the custom performance counter CSRs, all the implemented CSRs follow the standard mapping of bit fields as described in the RISC\-/V Instruction Set Manual, Volume II: Privileged Architecture, Version 1.10. It must be noted that even among the standard CSRs, not all bit fields have been implemented, limited by the subset of features implemented in the CPU. Refer to the next section for detailed description of the subset of fields implemented under each of these CSRs.

\subfile{\modulefiles/01-RISCV-CPU-reg-summ__EN}

Note that if write/set/clear operation is attempted on any of the CSRs which are read-only (RO), as indicated in the above table, the CPU will generate illegal instruction exception.

\subsection{Register Description}

\subfile{\modulefiles/01-RISCV-CPU-reg__EN}

%%%%%%%%%%%%%%%%%%%%%%%%%%%%
%%% Interrupt Controller %%%
%%%%%%%%%%%%%%%%%%%%%%%%%%%%

\clearpage
\section{Interrupt Controller}
\label{sec:riscvcpu-int-ctrl}

\subsection{Features}
\label{sec:riscvcpu-int-features}

The interrupt controller allows capturing, masking and dynamic prioritization of interrupt sources routed from peripherals to the RISC-V CPU. It has the following features:

\begin{itemize}
    \item Up to 31 asynchronous interrupts with unique IDs (1-31)
    \item Configurable via read/write to memory mapped registers
    \item 15 levels of priority, programmable for each interrupt
    \item Support for both level and edge type interrupt sources
    \item Programmable global threshold for masking interrupts with lower priority
    \item Interrupts IDs mapped to trap-vector address offsets
\end{itemize}


\subsection{Functional Description}
\label{sec:riscvcpu-int-func-descr}

Each interrupt ID has 5 properties associated with it:

\begin{enumerate}
    \item Enable State (0-1):
    \begin{itemize}
        \item Determines if an interrupt is enabled to be captured and serviced by the CPU.
        \item Programmed by writing the corresponding bit in \hyperref[regdesc:INTERRUPTCORE0CPUINTENABLEREG]{INTERRUPT\_CORE0\_CPU\_INT\_ENABLE\_REG}.
    \end{itemize}
    \item Type (0-1):
    \begin{itemize}
        \item Enables latching the state of an interrupt signal on its rising edge.
        \item Programmed by writing the corresponding bit in \hyperref[regdesc:INTERRUPTCORE0CPUINTTYPEREG]{INTERRUPT\_CORE0\_CPU\_INT\_TYPE\_REG}.
        \item An interrupt for which type is kept 0 is referred as a 'level' type interrupt.
        \item An interrupt for which type is set to 1 is referred as an 'edge' type interrupt.
    \end{itemize}
    \item Priority (1-15):
    \begin{itemize}
        \item Determines which interrupt, among multiple pending interrupts, the CPU will service first.
        \item Programmed by writing to the \hyperref[regdesc:INTERRUPTCORE0CPUINTPRI0REG]{INTERRUPT\_CORE0\_CPU\_INT\_PRI\_\regindex{n}\_REG} for a particular interrupt ID \regindex{n} in range (1-31).
        \item Enabled interrupts with priorities zero or less than the threshold value in \hyperref[regdesc:INTERRUPTCORE0CPUINTTHRESHREG]{INTERRUPT\_CORE0\_CPU\_INT\_THRESH\_REG} are masked.
        \item Priority levels increase from 1 (lowest) to 15 (highest).
        \item Interrupts with same priority are statically prioritized by their IDs, lowest ID having highest priority.
    \end{itemize}
    \item Pending State (0-1):
    \begin{itemize}
        \item Reflects the captured state of an enabled and unmasked interrupt signal.
        \item For each interrupt ID, the corresponding bit in read-only \hyperref[regdesc:INTERRUPTCORE0CPUINTEIPSTATUSREG]{INTERRUPT\_CORE0\_CPU\_INT\_EIP\_STATUS\_REG} gives its pending state.
        \item A pending interrupt will cause CPU to enter trap if no other pending interrupt has higher priority.
        \item A pending interrupt is said to be 'claimed' if it preempts the CPU and causes it to jump to the corresponding trap vector address.
        \item All pending interrupts which are yet to be serviced are termed as 'unclaimed'.
    \end{itemize}
    \item Clear State (0-1):
    \begin{itemize}
        \item Toggling this will clear the pending state of claimed edge-type interrupts only.
        \item Toggled by first setting and then clearing the corresponding bit in \hyperref[regdesc:INTERRUPTCORE0CPUINTCLEARREG]{INTERRUPT\_CORE0\_CPU\_INT\_CLEAR\_REG}.
        \item Pending state of a level type interrupt is unaffected by this and must be cleared from source.
        \item Pending state of an unclaimed edge type interrupt can be flushed, if required, by first clearing the corresponding bit in \hyperref[regdesc:INTERRUPTCORE0CPUINTENABLEREG]{INTERRUPT\_CORE0\_CPU\_INT\_ENABLE\_REG} and then toggling same bit in \hyperref[regdesc:INTERRUPTCORE0CPUINTCLEARREG]{INTERRUPT\_CORE0\_CPU\_INT\_CLEAR\_REG}.
    \end{itemize}
\end{enumerate}

When CPU services a pending interrupt, it:
\begin{itemize}
    \item saves the address of the current un-executed instruction in \hyperref[regdesc:MEPC]{mepc} for resuming execution later.
    \item updates the value of \hyperref[regdesc:MCAUSE]{mcause} with the ID of the interrupt being serviced.
    \item copies the state of \hyperref[fielddesc:MSTATUSMIE]{MIE} into \hyperref[fielddesc:MSTATUSMPIE]{MPIE}, and subsequently clears \hyperref[fielddesc:MSTATUSMIE]{MIE}, thereby disabling interrupts globally.
    \item enters trap by jumping to a word-aligned offset of the address stored in \hyperref[regdesc:MTVEC]{mtvec}.
\end{itemize}

Table \ref{tab:riscvcpu-int-isr-addr-map} shows the mapping of each interrupt ID with the corresponding trap-vector address. In short, the word aligned trap address for an interrupt with a certain $ID=i$ can be calculated as $(\hyperref[fielddesc:MTVECBASE]{mtvec} + 4i)$.

 \textit{Note: $ID=0$ is unavailable and therefore cannot be used for capturing interrupts. This is because the corresponding trap vector address (\hyperref[fielddesc:MTVECBASE]{mtvec} + 0x00) is reserved for exceptions.}

\begin{longtable}{ | r | l || r | l || r | l || r | l | }
\caption{ID wise map of Interrupt Trap-Vector Addresses} \label{tab:riscvcpu-int-isr-addr-map}
\\ \hline
\rowcolor{lightgray}
ID & Address & ID & Address & ID & Address & ID & Address \\ \hline
 0 & NA           &  8 & mtvec + 0x20 & 16 & mtvec + 0x40 & 24 & mtvec + 0x60   \\ \hline
 1 & mtvec + 0x04 &  9 & mtvec + 0x24 & 17 & mtvec + 0x44 & 25 & mtvec + 0x64   \\ \hline
 2 & mtvec + 0x08 & 10 & mtvec + 0x28 & 18 & mtvec + 0x48 & 26 & mtvec + 0x68   \\ \hline
 3 & mtvec + 0x0c & 11 & mtvec + 0x2c & 19 & mtvec + 0x4c & 27 & mtvec + 0x6c   \\ \hline
 4 & mtvec + 0x10 & 12 & mtvec + 0x30 & 20 & mtvec + 0x50 & 28 & mtvec + 0x70   \\ \hline
 5 & mtvec + 0x14 & 13 & mtvec + 0x34 & 21 & mtvec + 0x54 & 29 & mtvec + 0x74   \\ \hline
 6 & mtvec + 0x18 & 14 & mtvec + 0x38 & 22 & mtvec + 0x58 & 30 & mtvec + 0x78   \\ \hline
 7 & mtvec + 0x1c & 15 & mtvec + 0x3c & 23 & mtvec + 0x5c & 31 & mtvec + 0x7c   \\ \hline
\end{longtable}

After jumping to the trap-vector, the execution flow is dependent on software implementation, although it can be presumed that the interrupt will get handled (and cleared) in some interrupt service routine (ISR) and later the normal execution will resume once the CPU encounters MRET instruction.

Upon execution of MRET instruction, the CPU:
\begin{itemize}
    \item copies the state of \hyperref[fielddesc:MSTATUSMPIE]{MPIE} back into \hyperref[fielddesc:MSTATUSMIE]{MIE}, and subsequently clears \hyperref[fielddesc:MSTATUSMPIE]{MPIE}. This means that if previously \hyperref[fielddesc:MSTATUSMPIE]{MPIE} was set, then, after MRET, \hyperref[fielddesc:MSTATUSMIE]{MIE} will be set, thereby enabling interrupts globally.
    \item jumps to the address stored in \hyperref[regdesc:MEPC]{mepc} and resumes execution.
\end{itemize}

It is possible to perform software assisted nesting of interrupts inside an ISR as explained in \ref{sec:riscvcpu-int-suggest-op}.

The below listed points outline the functional behavior of the controller:

\begin{itemize}
    \item Only if an interrupt has non-zero priority, higher or equal to the value in the threshold register, will it be reflected in \hyperref[regdesc:INTERRUPTCORE0CPUINTEIPSTATUSREG]{INTERRUPT\_CORE0\_CPU\_INT\_EIP\_STATUS\_REG}.
    \item If an interrupt is visible in \hyperref[regdesc:INTERRUPTCORE0CPUINTEIPSTATUSREG]{INTERRUPT\_CORE0\_CPU\_INT\_EIP\_STATUS\_REG} and has yet to be serviced, then it's possible to mask it (and thereby prevent the CPU from servicing it) by either lowering the value of its priority or increasing the global threshold.
    \item If an interrupt, visible in \hyperref[regdesc:INTERRUPTCORE0CPUINTEIPSTATUSREG]{INTERRUPT\_CORE0\_CPU\_INT\_EIP\_STATUS\_REG}, is to be flushed (and prevented from being serviced at all), then it must be disabled (and cleared if it is of edge type).
\end{itemize}

\subsection{Suggested Operation}
\label{sec:riscvcpu-int-suggest-op}

\subsubsection{Latency Aspects}
\label{sec:riscvcpu-int-latency}

There is latency involved while configuring the Interrupt Controller.

In steady state operation, the Interrupt Controller has a fixed latency of 4 cycles. Steady state means that no changes have been made to the Interrupt Controller registers recently. This implies that any interrupt that is asserted to the controller will take exactly 4 cycles before the CPU starts processing the interrupt. This further implies that CPU may execute up to 5 instructions before the preemption happens.

Whenever any of its registers are modified, the Interrupt Controller enters into transient state, which may take up to 4 cycles for it to settle down into steady state again. During this transient state, the ordering of interrupts may not be predictable, and therefore, a few safety measures need to be taken in software to avoid any synchronization issues.

Also, it must be noted that the Interrupt Controller configuration registers lie in the APB address range, hence any R/W access to these registers may take multiple cycles to complete.

In consideration of above mentioned characteristics, users are advised to follow the sequence described below, whenever modifying any of the Interrupt Controller registers:
\begin{enumerate}
    \item save the state of \hyperref[fielddesc:MSTATUSMIE]{MIE} and clear \hyperref[fielddesc:MSTATUSMIE]{MIE} to 0
    \item read-modify-write one or more Interrupt Controller registers
    \item execute FENCE instruction to wait for any pending write operations to complete
    \item finally, restore the state of \hyperref[fielddesc:MSTATUSMIE]{MIE}
\end{enumerate}

Due to its critical nature, it is recommended to disable interrupts globally (\hyperref[fielddesc:MSTATUSMIE]{MIE}=0) beforehand, whenever configuring interrupt controller registers, and then restore \hyperref[fielddesc:MSTATUSMIE]{MIE} right after, as shown in the sequence above.

After execution of the sequence above, the Interrupt Controller will resume operation in steady state.

\subsubsection{Configuration Procedure}
\label{sec:riscvcpu-int-config}

By default, interrupts are disabled globally, since the reset value of \hyperref[fielddesc:MSTATUSMIE]{MIE} bit in \hyperref[regdesc:MSTATUS]{mstatus} is 0. Software must set \hyperref[fielddesc:MSTATUSMIE]{MIE}=1 after initialization of the interrupt stack (including setting \hyperref[regdesc:MTVEC]{mtvec} to the interrupt vector address) is done.

During normal execution, if an interrupt \regindex{n} is to be enabled, the below sequence may be followed:
\begin{enumerate}
    \item save the state of \hyperref[fielddesc:MSTATUSMIE]{MIE} and clear \hyperref[fielddesc:MSTATUSMIE]{MIE} to 0
    \item depending upon the type of the interrupt (edge/level), set/unset the \regindex{n}th bit of \hyperref[regdesc:INTERRUPTCORE0CPUINTTYPEREG]{INTERRUPT\_CORE0\_CPU\_INT\_TYPE\_REG}
    \item set the priority by writing a value to \hyperref[regdesc:INTERRUPTCORE0CPUINTPRI0REG]{INTERRUPT\_CORE0\_CPU\_INT\_PRI\_\regindex{n}\_REG} in range 1 (lowest) to 15 (highest)
    \item set the \regindex{n}th bit of \hyperref[regdesc:INTERRUPTCORE0CPUINTENABLEREG]{INTERRUPT\_CORE0\_CPU\_INT\_ENABLE\_REG}
    \item execute FENCE instruction
    \item restore the state of \hyperref[fielddesc:MSTATUSMIE]{MIE}
\end{enumerate}

When one or more interrupts become pending, the CPU acknowledges (claims) the interrupt with the highest priority and jumps to the trap vector address corresponding to the interrupt's ID. Software implementation may read \hyperref[regdesc:MCAUSE]{mcause} to infer the type of trap (\hyperref[fielddesc:MCAUSEINT]{mcause(31)} is 1 for interrupts and 0 for exceptions) and then the ID of the interrupt (\hyperref[fielddesc:MCAUSECODE]{mcause(4-0)} gives ID of interrupt or exception). This inference may not be necessary if each entry in the trap vector are jump instructions to different trap handlers. Ultimately, the trap handler(s) will redirect execution to the appropriate ISR for this interrupt.

Upon entering into an ISR, software must toggle the \regindex{n}th bit of \hyperref[regdesc:INTERRUPTCORE0CPUINTCLEARREG]{INTERRUPT\_CORE0\_CPU\_INT\_CLEAR\_REG} if the interrupt is of edge type, or clear the source of the interrupt if it is of level type.

Software may also update the value of \hyperref[regdesc:INTERRUPTCORE0CPUINTTHRESHREG]{INTERRUPT\_CORE0\_CPU\_INT\_THRESH\_REG} and program \hyperref[fielddesc:MSTATUSMIE]{MIE}=1 for allowing higher priority interrupts to preempt the current ISR (nesting), however, before doing so, all the state CSRs must be saved (\hyperref[regdesc:MEPC]{mepc}, \hyperref[regdesc:MSTATUS]{mstatus}, \hyperref[regdesc:MCAUSE]{mcause}, etc.) since they will get overwritten due to occurrence of such an interrupt. Later, when exiting the ISR, the values of these CSRs must be restored.

Finally, after the execution returns from the ISR back to the trap handler, MRET instruction is used to resume normal execution.

Later, if the \regindex{n} interrupt is no longer needed and needs to be disabled, the following sequence may be followed:
\begin{enumerate}
    \item save the state of \hyperref[fielddesc:MSTATUSMIE]{MIE} and clear \hyperref[fielddesc:MSTATUSMIE]{MIE} to 0
    \item check if the interrupt is pending in \hyperref[regdesc:INTERRUPTCORE0CPUINTEIPSTATUSREG]{INTERRUPT\_CORE0\_CPU\_INT\_EIP\_STATUS\_REG}
    \item set/unset the \regindex{n}th bit of \hyperref[regdesc:INTERRUPTCORE0CPUINTENABLEREG]{INTERRUPT\_CORE0\_CPU\_INT\_ENABLE\_REG}
    \item if the interrupt is of edge type and was found to be pending in step 2 above, \regindex{n}th bit of \hyperref[regdesc:INTERRUPTCORE0CPUINTCLEARREG]{INTERRUPT\_CORE0\_CPU\_INT\_CLEAR\_REG} must be toggled, so that its pending status gets flushed
    \item execute FENCE instruction
    \item restore the state of \hyperref[fielddesc:MSTATUSMIE]{MIE}
\end{enumerate}

Above is only a suggested scheme of operation. Actual software implementation may vary.

%%%%%%%%%%%%%%%%%%%%%%%%
%%% Register Summary %%%
%%%%%%%%%%%%%%%%%%%%%%%%

\hypertarget{riscvcpu-int-reg-summ}{}
\subsection{Register Summary}
\label{sec:riscvcpu-int-reg-summ}


The addresses in this section are relative to Interrupt Controller base address provided in Table \ref{tab:sysmem-base-address} in Chapter \ref{mod:sysmem} \textit{\nameref{mod:sysmem}}.

For the complete list of interrupt registers and detailed configuration information, please refer to Chapter \ref{mod:intmtrx} \textit{\nameref{mod:intmtrx}}, section \ref{sec:intmtrx-reg-summ}, register group "CPU Interrupt Registers".


%%%%%%%%%%%%%%%%%
%%% Registers %%%
%%%%%%%%%%%%%%%%%

\subsection{Register Description}

The addresses in this section are relative to Interrupt Controller base address provided in Table \ref{tab:sysmem-base-address} in Chapter \ref{mod:sysmem} \textit{\nameref{mod:sysmem}}.

For the complete list of interrupt registers and detailed configuration information, please refer to Chapter \ref{mod:intmtrx} \textit{\nameref{mod:intmtrx}}, section \ref{sec:intmtrx-reg-summ}, register group "CPU interrupt register".



%%%%%%%%%%%%%%%%%%%%%%%%%%%%
%%% Debug %%%
%%%%%%%%%%%%%%%%%%%%%%%%%%%%

\clearpage
\section{Debug}
\label{sec:riscvcpu-dbg}


\subsection{Overview}
\label{sec:riscvcpu-dbg-overview}
This section describes how to debug and test software running on CPU core. Debug support is provided through standard JTAG pins and complies to RISC-V External Debug Support Specification version 0.13.

Figure \ref{fig:riscvcpu-dbg-sys-overview} below shows the main components of External Debug Support.

\begin{figure}[h]
\centering
\includegraphics[width=\textwidth/2,keepaspectratio]{\modulefiles/figures/debug_bigfonts_22feb2022.png}
\caption{Debug System Overview}\label{fig:riscvcpu-dbg-sys-overview}
\end{figure}

The user interacts with the Debug Host (e.g., laptop), which is running a debugger (e.g., gdb). The debugger communicates with a Debug Translator (e.g., OpenOCD, which may include a hardware
driver) to communicate with Debug Transport Hardware (e.g., Olimex USB-JTAG adapter). The Debug Transport Hardware connects the Debug Host to the ESP-RV Core's Debug Transport Module (DTM) through standard JTAG interface. The DTM provides access to the Debug Module (DM) using the Debug Module Interface (DMI).

The DM allows the debugger to halt the core. Abstract commands provide access to its GPRs (general purpose registers). The Program Buffer allows the debugger to execute arbitrary code on the core, which allows access to additional CPU core state. Alternatively, additional abstract commands can provide access to additional CPU core state.
ESP-RV core contains Trigger Module supporting two triggers. When trigger conditions are met, cores will halt spontaneously and inform the debug module that they have halted.

System bus access block allows memory and peripheral register access without using RISC-V core.


\subsection{Features}
\label{sec:riscvcpu-dbg-features}
Basic debug functionality has the following features:

\begin{itemize}
    \item {Halt and resume CPU core}
    \item {Access to CSR and GPR}
    \item {Debug from the first instruction after reset}
    \item {Core reset control}
    \item {Software breakpoint}
    \item {Hardware single-stepping}
    \item {16-word program buffer}
    \item {System bus access}
    \item {Support for two hardware triggers}
\end{itemize}

\subsection{Functional Description}
\label{sec:riscvcpu-dbg-func-descr}

As mentioned earlier, Debug Scheme conforms to RISC-V External Debug Support Specification version 0.13. Please refer the specs for functional operation details.

%%%%%%%%%%%%%%%%%%%%%%%%
%%% Register Summary %%%
%%%%%%%%%%%%%%%%%%%%%%%%

\hypertarget{riscvcpu-dbg-reg-summ}{}
\subsection{Register Summary}
\label{sec:riscvcpu-dbg-reg-summ}


Below is the list of Debug CSR's supported by ESP-RV core.
\subfile{\modulefiles/01-RISCV-DBG-reg-summ__EN}

All the debug module registers are implemented in conformance to RISC-V External Debug Support Specification version 0.13. Please refer it for more details.

%%%%%%%%%%%%%%%%%
%%% Registers %%%
%%%%%%%%%%%%%%%%%


\subsection{Register Description}
Below are the details of Debug CSR's supported by ESP-RV core.

\subfile{\modulefiles/01-RISCV-DBG-reg__EN}

%%%%%%%%%%%%%%%%%%%%%%%%
%%% Hardware Trigger %%%
%%%%%%%%%%%%%%%%%%%%%%%%

\clearpage
\section{Hardware Trigger}
\label{sec:riscvcpu-hwt}


\subsection{Features}
\label{sec:riscvcpu-hwt-features}

Hardware Trigger module provides breakpoint and watchpoint capability for debugging. It has the following features:

\begin{itemize}
    \item Two independent trigger units
    \item Matching the address of program counter or load-store accesses
    \item Execution preemption by causing breakpoint exception
    \item Halting execution and transferring control to debugger
    \item Support for NAPOT (naturally aligned power of two) address encoding
\end{itemize}

\subsection{Functional Description}
\label{sec:riscvcpu-hwt-func-descr}

The Hardware Trigger module provides four CSRs, which are listed under \hyperref[sec:riscvcpu-hwt-reg-summ]{register summary} section. Among these, \hyperref[regdesc:TDATA1]{tdata1} and \hyperref[regdesc:TDATA2]{tdata2} are abstract CSRs, which means they are shadow registers for accessing internal registers for each of the eight trigger units, one at a time.

To choose a particular trigger unit write the index (0-7) of that unit into \hyperref[regdesc:TSELECT]{tselect} CSR. When tselect is written with a valid index, the abstract CSRs tdata1 and tdata2 are automatically mapped to reflect internal registers of that trigger unit. Each trigger unit has two internal registers, namely \hyperref[regdesc:MCONTROL]{mcontrol} and \hyperref[regdesc:MADDRESS]{maddress}, which are mapped to tdata1 and tdata2, respectively.

Writing larger than allowed indexes to tselect will clip the written value to the largest valid index, which can be read back. This property may be used for enumerating the number of available triggers during initialization or when using a debugger.

Since software or debugger may need to know the type of the selected trigger to correctly interpret tdata1 and tdata2, the 4 bits
(31-28) of tdata1 encodes the type of the selected trigger. This type field is read-only and always provides a value of 0x2 for every trigger, which stands for match type trigger, hence, it is inferred that tdata1 and tdata2 are to be interpreted as mcontrol and maddress. The information regarding other possible values can be found in the RISC-V Debug Specification v0.13, but this trigger module only supports type 0x2.

Once a trigger unit has been chosen by writing its index to tselect, it will become possible to configure it by setting the appropriate bits in mcontrol CSR (tdata1) and writing the target address to maddress CSR (tdata2).

Each trigger unit can be configured to either cause breakpoint exception or enter debug mode, by writing to the action bit of mcontrol. This bit can only be written from debugger, thus by default a trigger, if enabled, will cause breakpoint exception.

mcontrol for each trigger unit has a \hyperref[fielddesc:MCONTROLHIT]{hit} bit which may be read, after CPU halts or enters exception, to find out if this was the trigger unit that fired. This bit is set as soon as the corresponding trigger fires, but it has to be manually cleared before resuming operation. Although, failing to clear it doesn't affect normal execution in any way.

Each trigger unit only supports match on address, although this address could either be that of a load/store access or the virtual address of an instruction. The address and size of a region are specified by writing to maddress (tdata2) CSR for the selected trigger unit. Larger than 1 byte region sizes are specified through NAPOT (naturally aligned power of two) encoding (see Table \ref{tab:riscvcpu-hwt-table}) and enabled by setting match bit in mcontrol. Note that for NAPOT encoded addresses, by definition, the start address is constrained to be aligned to (i.e., an integer multiple of) the region size.

\begin{longtable}{ | c | c | r | }
\caption{NAPOT encoding for maddress} \label{tab:riscvcpu-hwt-table}
\\ \hline
\rowcolor{lightgray}
maddress (31-0) & Start Address & Size (bytes) \\ \hline
aaa...aaaaaaaaa0 & aaa...aaaaaaaaa0 & $2$ \\ \hline
aaa...aaaaaaaa01 & aaa...aaaaaaaa00 & $4$ \\ \hline
aaa...aaaaaaa011 & aaa...aaaaaaa000 & $8$ \\ \hline
aaa...aaaaaa0111 & aaa...aaaaaa0000 & $16$ \\ \hline
\multicolumn{3}{c}{....} \\ \hline
a01...1111111111 & a00...0000000000 & $2^{31}$ \\ \hline
\end{longtable}

\hyperref[regdesc:TCONTROL]{tcontrol} CSR is common to all trigger units. It is used for preventing triggers from causing repeated exceptions in machine-mode while execution is happening inside a trap handler. This also disables breakpoint exceptions inside ISRs by default, although, it is possible to manually enable this right before entering an ISR, for debugging purposes. This CSR is not relevant if a trigger is configured to enter debug mode.

\subsection{Trigger Execution Flow}
\label{sec:riscvcpu-hwt-suggest-op}

When hart is halted and enters debug mode due to the firing of a trigger (\hyperref[fielddesc:MCONTROLACTION]{action} = 1):
\begin{itemize}
    \item \hyperref[regdesc:DPC]{dpc} is set to current PC (in decode stage)
    \item cause field in \hyperref[regdesc:DCSR]{dcsr} is set to 2, which means halt due to trigger
    \item \hyperref[fielddesc:MCONTROLHIT]{hit} bit is set to 1, corresponding to the trigger(s) which fired
\end{itemize}

When hart goes into trap due to the firing of a trigger (\hyperref[fielddesc:MCONTROLACTION]{action} = 0) :
\begin{itemize}
    \item \hyperref[regdesc:MEPC]{mepc} is set to current PC (in decode stage)
    \item \hyperref[regdesc:MCAUSE]{mcause} is set to 3, which means breakpoint exception
    \item \hyperref[fielddesc:TCONTROLMPTE]{mpte} is set to the value in \hyperref[fielddesc:TCONTROLMTE]{mte} right before trap
    \item \hyperref[fielddesc:TCONTROLMTE]{mte} is set to 0
    \item \hyperref[fielddesc:MCONTROLHIT]{hit} bit is set to 1, corresponding to the trigger(s) which fired
\end{itemize}

\textit{Note : If two different triggers fire at the same time, one with action = 0 and another with action = 1, then hart is halted and enters debug mode.}

%%%%%%%%%%%%%%%%%%%%%%%%
%%% Register Summary %%%
%%%%%%%%%%%%%%%%%%%%%%%%

\hypertarget{riscvcpu-hwt-reg-summ}{}
\subsection{Register Summary}
\label{sec:riscvcpu-hwt-reg-summ}


Below is a list of Trigger Module CSRs supported by the CPU. These are only accessible from machine-mode.

\subfile{\modulefiles/01-RISCV-HWT-reg-summ__EN}

%%%%%%%%%%%%%%%%%
%%% Registers %%%
%%%%%%%%%%%%%%%%%

\subsection{Register Description}

\subfile{\modulefiles/01-RISCV-HWT-reg__EN}

%%%%%%%%%%%%%%%%%%%%%%%%
%%%%%%%%%%%%%%%%%%%%%%%%

\clearpage
\section{Memory Protection}
\label{sec:riscvcpu-pmp}

\subsection{Overview}
\label{sec:riscvcpu-pmp-overview}
The CPU core includes a physical memory protection unit, which can be used by software to set memory access privileges (read, write and execute permissions) for required memory regions. It supports 16 memory regions, of which some address regions have been hard coded to values in accordance with \chipname{} memory map and the rest of address regions are kept programmable to split SRAM into separate IRAM/DRAM regions as per software code size.

It is fully compliant to the  Physical Memory Protection (PMP) description specified in RISC\-/V Instruction Set Manual, Volume II: Privileged Architecture, Version 1.10. However, in order to save area, values of 13 \texttt{pmpaddrX} registers (refer \hyperref[sec:riscvcpu-pmp-reg-summ]{Register Summary}) have been hard-coded. Details are provided in next sub-section.

For detailed understanding of the RISC-V PMP concept, please refer to RISC\-/V Instruction Set Manual, Volume II: Privileged Architecture, Version 1.10.

\subsection{Features}
\label{sec:riscvcpu-pmp-features}
The memory protection unit has the following features:
\begin{itemize}
    \item {Support for 16 PMP entries}
    \item {Programmable \texttt{pmpaddr0-2}}
    \item {Hard-coded \texttt{pmpaddr3-15} as per \chipname{} memory map}
\end{itemize}

\subsection{Functional Description}
\label{sec:riscvcpu-pmp-func-descr}
Software can program the PMP unit's configuration and address registers in order to contain faults and support secure execution. PMP CSR's can only be programmed in machine-mode. Once enabled, write, read and execute permission checks are applied to all the accesses in user-mode as per programmed values of enabled \texttt{pmpcfgX} and \texttt{pmpaddrX} registers.

By default, PMP grants permission to all accesses in machine-mode and revokes permission of all access in user-mode. This implies that it is mandatory to program address range and valid permissions in \texttt{pmpcfgX} and \texttt{pmpaddr} registers for any valid access to pass through in user-mode. However, it is not required for machine-mode as PMP permits all accesses to go through by default. In cases where PMP checks are also required in machine-mode, software can set the lock bit of required PMP entry to enable permission checks on it. Once lock bit is set, it can only be cleared through CPU reset.

When any instruction is being fetched from memory region without execute permissions, exception is generated at processor level and exception cause is set as instruction access fault in \hyperref[regdesc:MCAUSE]{mcause} CSR. Similarly, any load/store access without valid read/write permissions, will result in exception generation with \hyperref[regdesc:MCAUSE]{mcause} updated as load access and store access fault respectively. In case of load/store access faults, violating address is captured in \hyperref[regdesc:MTVAL]{mtval} CSR.

%%%%%%%%%%%%%%%%%%%%%%%%
%%% Register Summary %%%
%%%%%%%%%%%%%%%%%%%%%%%%

\hypertarget{riscvcpu-pmp-reg-summ}{}
\subsection{Register Summary}
\label{sec:riscvcpu-pmp-reg-summ}


Below is a list of PMP CSRs supported by the CPU. These are only accessible from machine-mode. As mentioned earlier, \texttt{pmpaddrX0-2} are kept programmable to split the SRAM region as per software requirements. \texttt{pmpaddrX3-15} are hard coded with values specified in "CSR Reset Value" column.  These hard-coded values have been derived to match the SoC memory map specified in "PMP Region" column. To enable any PMP region, the A field in associated \texttt{pmpcfgX} register should always be programmed with value specified in "Address Matching Mode" column.

\subfile{\modulefiles/01-RISCV-PMP-reg-summ__EN}

%%%%%%%%%%%%%%%%%
%%% Registers %%%
%%%%%%%%%%%%%%%%%

\subsection{Register Description}
\label{sec:riscvcpu-pmp-reg-desc}

PMP unit implements all \texttt{pmpcfg0-3} and \texttt{pmpaddr0-15} CSRs as defined in RISC-V Instruction Set Manual Volume II: Privileged Architecture, Version 1.10.
\end{document}
